% Generated draft: requires mathematical and source review.
\documentclass[11pt]{article}
\usepackage[margin=25mm]{geometry}
\usepackage{fontspec}
\setmainfont{Latin Modern Roman}
\usepackage{amsmath,amssymb,mathtools,unicode-math}
\setmathfont{Latin Modern Math}
\usepackage{enumitem,needspace,xurl,hyperref}
\setlength{\parindent}{0pt}
\setlength{\parskip}{5pt}
\setlength{\emergencystretch}{3em}
\newcommand{\N}{\mathbb{N}}
\newcommand{\Z}{\mathbb{Z}}
\newcommand{\Q}{\mathbb{Q}}
\newcommand{\R}{\mathbb{R}}
\newcommand{\C}{\mathbb{C}}
\newcommand{\F}{\mathbb{F}}
\newcommand{\A}{\mathbb{A}}
\newcommand{\PP}{\mathbb{P}}
\newcommand{\E}{\mathbb{E}}
\newcommand{\eps}{\varepsilon}
\newcommand{\dd}{\,\mathrm{d}}
\begin{document}
\section*{UnsolvedMath ID: 800248}
\textbf{Source problem code:} catalogue entry
\subsection{Definitions and mathematical statement}
Let \(X\subset\PP^n(\C)\) be a compact complex hypersurface of degree \(d\ge n+2\), having only normal crossing singularities, and put \(M=\PP^n(\C)\setminus X\). The source does not explicitly specify a range for \(n\); the natural setting is a positive integer \(n\). Under the usual reduced local interpretation of normal crossings, near each point of \(X\) there are holomorphic coordinates in which \(X\) has equation \(z_1\cdots z_r=0\) for some \(1\le r\le n\). The source does not explicitly specify reducedness or whether it intends a stronger global simple normal crossing condition.

Write \(\Delta=\{\zeta\in\C:|\zeta|<1\}\). For \(p\in M\) and \(v\in T_pM\), define the infinitesimal Kobayashi pseudometric by
\[
\kappa_M(p;v)=\inf\left\{\frac{1}{a}:a>0,\ f:\Delta\to M\text{ holomorphic},\ f(0)=p,\ \dd f_0\!\left(\frac{\partial}{\partial\zeta}\right)=av\right\}.
\]
Its degeneracy locus is
\[
D_M=\{p\in M:\text{there exists }v\in T_pM\setminus\{0\}\text{ with }\kappa_M(p;v)=0\}.
\]
\textbf{Requested assertion, with an explicit interpretation of the target.} Show that there exists a complex analytic hypersurface \(H\subset M\), allowing \(H=\varnothing\), such that
\[
D_M\subseteq H.
\]
Equivalently, \(\kappa_M(p;v)>0\) for every \(p\in M\setminus H\) and every nonzero \(v\in T_pM\). Here a nonempty complex analytic hypersurface means a closed analytic subset of pure complex codimension one in \(M\). This interprets the source phrase ``degenerate at most on a complex hypersurface'' as containment of the degeneracy locus in such a hypersurface. The source does not specify whether the exceptional hypersurface must instead be algebraic, extend to a hypersurface of \(\PP^n(\C)\), or satisfy another global condition; that qualification remains unresolved.
\subsection{Short English statement}
For a compact hypersurface \(X\subset\PP^n(\C)\) of degree at least \(n+2\) with only normal crossing singularities, prove that the infinitesimal Kobayashi pseudometric of its complement is positive on every nonzero tangent vector outside an exceptional complex hypersurface. The precise normal crossing convention and the intended global meaning of the exceptional hypersurface require review.
\subsection{Sources}
\begin{itemize}
\item Local UnsolvedMath registry: permanent ID \(800248\), source problem code \texttt{catalogue entry}, title \emph{Several Complex Variables Problem List --- Problem 47}. The supplied registry labels this record \texttt{open}; that label is supplied source metadata, not an independently verified mathematical status.
\item Buzzard, Fornaess, and Krantz, \emph{The Several Complex Variables Problem List} (1995), Problem 47. Supplied source URL: \url{https://arxiv.org/abs/math/9509202}. The registry lists Buzzard, Fornaess, and Krantz as proposers and gives the year 1995; the extracted problem itself is attributed to N. Sibony, with contact \texttt{sibony@anh.matups.fr}.
\item Supplied provenance: accessed 2026-10-16; extraction method \texttt{arxiv\_latex}; upstream revision \texttt{d371ee917ac9647fbb06cd6fc529deec6f3dda82}. This access date and provenance are supplied curation information; no independent source visit or literature verification is asserted.
\end{itemize}
\end{document}
